\section{Introduction}
\label{sec:intro}

Two prompts can look the same to a person and still differ in bytes.
A review pasted from a web page carries no-break spaces (U+00A0) instead of ordinary spaces.
A file saved on Windows ends its lines with \texttt{\textbackslash r\textbackslash n}.
A word processor replaces straight quotes with curly ones and hyphens with dashes.
None of these edits changes the meaning of the text, and most are invisible on screen.
Our main question is simple: {\bf do these invisible edits change what a language model outputs?}

The answer matters for anyone who measures or deploys LLMs.
If invisible formatting changes outputs, then benchmark numbers contain a hidden formatting term, A/B tests of prompt variants are confounded, and the same input behaves differently depending on the editor or operating system that produced it.
Small on-device models are a particular concern, because practitioners often assume they are brittle but rarely know by how much.

Prior work shows that LLMs are sensitive to prompt design.
Changing separators, casing or field names in a few-shot template can move accuracy by tens of points~\cite{sclar2024quantifying}, single-prompt evaluations give unstable rankings~\cite{mizrahi2024state}, and sensitivity varies across models and tasks~\cite{zhuo2024prosa}.
However, these studies perturb \emph{visible} template structure.
Character-level robustness work studies typos and adversarial edits~\cite{belinkov2018synthetic,pruthi2019combating,wang2021textflint}, which a reader can also see.
We lack a controlled comparison of edits that a reader cannot see, one that separates whitespace from punctuation, spans model sizes within families, covers both classification and code, and measures each effect against the noise floor of inference itself.

We fill this gap with a controlled study.
We define five whitespace edits (trailing spaces, CRLF line endings, double spaces, trailing newlines, and no-break spaces or tabs) and five punctuation edits (curly quotes, curly apostrophes, Unicode dashes, the ellipsis character, and all four combined).
We apply each edit to every item of \sst, \agnews, \boolq, \humaneval and \mbpp, and we run 11 open instruction-tuned models from the \qwentwo, \qwenthree and \llama families, ranging from 0.5B to 14.8B parameters.
Every perturbed output is compared with the same model's output on the canonical prompt.
Every item also gets a \emph{control} condition: the byte-identical canonical prompt placed at another batch position, which measures how often outputs change for no reason at all.
Finally, we pass every prompt through each model's chat template and tokenizer to check whether the edit survives and how many tokens it changes.

\Figref{fig:per_condition} summarizes the main result.
A single invisible edit flips 4.9\% of classification labels on average, more than 200 times the 0.02\% control rate.
We expected punctuation to matter more than whitespace. The opposite holds: whitespace edits flip 7.1\% of labels and punctuation edits 2.2\%, and whitespace is worse for all 11 models.
Part of the gap comes from token footprint, because a double-space or no-break-space edit re-tokenizes every word. Even at a matched footprint, though, whitespace is worse: a trailing \texttt{\textbackslash n\textbackslash n} (about one changed token) flips 3.4\% of labels, while a Unicode dash (1.3 tokens) flips 1.3\%.
Models under 3B parameters flip 2.2$\times$ as often as larger models on classification, but code shows no size effect.
Code text changes in 36\% of generations, yet only 5.1\% of programs change their test outcome.

Our contributions are:
\begin{itemize}[leftmargin=*,itemsep=0pt,topsep=0pt]
    \item We propose a taxonomy of ten invisible formatting edits, split into whitespace and punctuation families, with code-safe variants that leave programs valid.
    \item We conduct a controlled study over 11 open models, 5 tasks and 239,616 greedy generations, with a byte-identical re-run as the noise floor and an exact replication in a separate engine.
    \item We show that whitespace edits are 2--5$\times$ more disruptive than punctuation edits, that small models are about 2$\times$ more label-sensitive, and that most code changes are stylistic rather than functional.
    \item We trace the effect to the token level. Edits that the chat template removes, such as trailing newlines under \llama templates, have zero effect, and divergence grows with the number of changed tokens without fully explaining the whitespace--punctuation gap.
\end{itemize}

\Secref{sec:related} reviews related work, \secref{sec:method} describes our setup, \secref{sec:results} reports results, and \secref{sec:discussion} discusses mechanisms, implications and limitations.
